\section{Билет 14. Затухающие колебания. Уравнение затухающих колебаний и его решение. Частота затухающих колебаний. Зависимость амплитуды затухающих колебаний от времени. Графики затухающих колебаний при разных значениях коэффициента затухания. Время затухания.}

\begin{definition}
При условии малого затухания ($\beta < \omega_0$) решением уравнения \eqref{eq:damped_ode} является закон движения:
\begin{equation}
    x(t) = A_0 e^{-\beta t} \cos(\omega t + \alpha), \label{eq:damped_solution}
\end{equation}
где $A_0$ --- начальная амплитуда, $\alpha$ --- начальная фаза, определяемые из начальных условий, а $\omega$ --- частота затухающих колебаний.
\end{definition}

\begin{property}[Частота и период затухающих колебаний]
Циклическая частота затухающих колебаний равна:
\begin{equation}
    \omega = \sqrt{\omega_0^2 - \beta^2}. \label{eq:damped_freq}
\end{equation}
Так как $\beta > 0$, частота затухающих колебаний всегда меньше собственной: $\omega < \omega_0$.
Период затухающих колебаний определяется как:
\begin{equation}
    T = \frac{2\pi}{\omega} = \frac{2\pi}{\sqrt{\omega_0^2 - \beta^2}}. \label{eq:damped_period}
\end{equation}
Отсюда следует, что $T > T_0$ (период затухающих колебаний больше периода собственных). При незначительном сопротивлении среды ($\beta \ll \omega_0$) период практически совпадает с периодом незатухающих колебаний: $T \approx T_0 = 2\pi/\omega_0$. Влияние небольшого трения на период гораздо меньше, чем на амплитуду.
\end{property}

\begin{definition}
Из решения \eqref{eq:damped_solution} следует, что \textbf{амплитуда затухающих колебаний} экспоненциально убывает со временем:
\begin{equation}
    A(t) = A_0 e^{-\beta t}. \label{eq:amplitude_decay}
\end{equation}
Скорость убывания амплитуды определяется величиной $\beta$, которую называют \textbf{коэффициентом затухания}.
\end{definition}

\begin{remark}
\textbf{Графики затухающих колебаний.} График $x(t)$ представляет собой колебательную кривую, заключённую между двумя экспоненциальными огибающими $y = \pm A_0 e^{-\beta t}$. При увеличении коэффициента затухания $\beta$:
\begin{enumerate}
    \item Амплитуда убывает быстрее (график «сжимается» к оси времени).
    \item Период колебаний увеличивается (расстояние между соседними максимумами растёт).
    \item При $\beta \to \omega_0$ колебательный режим сменяется апериодическим (система медленно возвращается в равновесие без колебаний).
\end{enumerate}
\end{remark}

\begin{theorem}[Время затухания]
Важной характеристикой скорости затухания является \textbf{время затухания (релаксации)} $\tau$ --- промежуток времени, за который амплитуда колебаний уменьшается в $e$ раз ($e \approx 2.718$).
Из условия $A(\tau) = A_0 / e$ и формулы \eqref{eq:amplitude_decay} получаем:
\begin{equation}
    A_0 e^{-\beta \tau} = \frac{A_0}{e} \quad \Rightarrow \quad e^{-\beta \tau} = e^{-1} \quad \Rightarrow \quad \beta \tau = 1.
\end{equation}
Отсюда:
\begin{equation}
    \tau = \frac{1}{\beta}. \label{eq:relaxation_time}
\end{equation}
Таким образом, коэффициент затухания $\beta$ обратен времени релаксации: $\beta = 1/\tau$.
\end{theorem}
